% GENERATED FILE. Edit canonical lessons, then run npm run generate:books.
% canonical-source-hash: fnv1a64:da29f7557fdb62b0
% canonical-lessons: TE-C145-nercuko, TE-C145-nerpu, TE-C145-maracipo, TE-C145-duku, TE-C145-digu

\chapter{To learn, To teach, To forget, To jump, To get down}
\label{ch:te-to-learn-to-teach-to-forget-to-jump-to-get-down}
\hlchaptermodality{145}

\begin{chapteropening}
\textbf{By the end of this chapter:} \emph{I can say to learn, to teach, to forget, to jump and to get down.}

Complete the last lesson of chapter 145: i can say to learn, to teach, to forget, to jump and to get down.
\end{chapteropening}

\section[nērcukō]{\te{నేర్చుకో} (nērcukō) — to learn}

\label{lesson:TE-C145-nercuko}



\begin{quote}
Before the new one: say the Telugu for to send, then the Telugu for to look for.
\end{quote}

\subsection*{You'll want to know: \te{నేర్చుకో}}
\textbf{\te{నేర్చుకో}} — \emph{nērcukō} — \textquotedblleft{}to learn\textquotedblright{}.

\begin{grammarlens}[title={What you've built}]
One more word for this chapter.
\end{grammarlens}

\subsection*{Your turn}
\begin{itemize}
\raggedright
  \item \emph{Say it:} \emph{nērcukō}
  \item \emph{Say it:} \emph{nērcukō}, once more
  \item \emph{Say it:} say \emph{nērcukō}
  \item \emph{Recall:} say the Telugu for to send, then the Telugu for to look for, then say \emph{nērcukō} again
\end{itemize}

\begin{tcolorbox}[breakable,colback=teal!4,colframe=teal!35!black,title={Before you move on}]
What does \te{నేర్చుకో} mean? (\textquotedblleft{}To learn\textquotedblright{}.) Say it once more.
\end{tcolorbox}
\section[nērpu]{\te{నేర్పు} (nērpu) — to teach}

\label{lesson:TE-C145-nerpu}



\begin{quote}
Before the new one: say the Telugu for to look for, then the Telugu for to learn.
\end{quote}

\subsection*{You'll want to know: \te{నేర్పు}}
\textbf{\te{నేర్పు}} — \emph{nērpu} — \textquotedblleft{}to teach\textquotedblright{}.

\begin{grammarlens}[title={What you've built}]
One more word for this chapter.
\end{grammarlens}

\subsection*{Your turn}
\begin{itemize}
\raggedright
  \item \emph{Say it:} \emph{nērpu}
  \item \emph{Say it:} \emph{nērpu}, once more
  \item \emph{Say it:} say \emph{nērpu}
  \item \emph{Recall:} say the Telugu for to look for, then the Telugu for to learn, then say \emph{nērpu} again
\end{itemize}

\begin{tcolorbox}[breakable,colback=teal!4,colframe=teal!35!black,title={Before you move on}]
What does \te{నేర్పు} mean? (\textquotedblleft{}To teach\textquotedblright{}.) Say it once more.
\end{tcolorbox}
\section[maracipō]{\te{మరచిపో} (maracipō) — to forget}

\label{lesson:TE-C145-maracipo}



\begin{quote}
Before the new one: say the Telugu for to learn, then the Telugu for to teach.
\end{quote}

\subsection*{You'll want to know: \te{మరచిపో}}
\textbf{\te{మరచిపో}} — \emph{maracipō} — \textquotedblleft{}to forget\textquotedblright{}.

\begin{grammarlens}[title={What you've built}]
One more word for this chapter.
\end{grammarlens}

\subsection*{Your turn}
\begin{itemize}
\raggedright
  \item \emph{Say it:} \emph{maracipō}
  \item \emph{Say it:} \emph{maracipō}, once more
  \item \emph{Say it:} say \emph{maracipō}
  \item \emph{Recall:} say the Telugu for to learn, then the Telugu for to teach, then say \emph{maracipō} again
\end{itemize}

\begin{tcolorbox}[breakable,colback=teal!4,colframe=teal!35!black,title={Before you move on}]
What does \te{మరచిపో} mean? (\textquotedblleft{}To forget\textquotedblright{}.) Say it once more.
\end{tcolorbox}
\section[dūku]{\te{దూకు} (dūku) — to jump}

\label{lesson:TE-C145-duku}



\begin{quote}
Before the new one: say the Telugu for to teach, then the Telugu for to forget.
\end{quote}

\subsection*{You'll want to know: \te{దూకు}}
\textbf{\te{దూకు}} — \emph{dūku} — \textquotedblleft{}to jump\textquotedblright{}.

\begin{grammarlens}[title={What you've built}]
One more word for this chapter.
\end{grammarlens}

\subsection*{Your turn}
\begin{itemize}
\raggedright
  \item \emph{Say it:} \emph{dūku}
  \item \emph{Say it:} \emph{dūku}, once more
  \item \emph{Say it:} say \emph{dūku}
  \item \emph{Recall:} say the Telugu for to teach, then the Telugu for to forget, then say \emph{dūku} again
\end{itemize}

\begin{tcolorbox}[breakable,colback=teal!4,colframe=teal!35!black,title={Before you move on}]
What does \te{దూకు} mean? (\textquotedblleft{}To jump\textquotedblright{}.) Say it once more.
\end{tcolorbox}
\section[digu]{\te{దిగు} (digu) — to get down, to get off}

\label{lesson:TE-C145-digu}



\begin{quote}
Before the new one: say the Telugu for to forget, then the Telugu for to jump.
\end{quote}

\subsection*{You'll want to know: \te{దిగు}}
\textbf{\te{దిగు}} — \emph{digu} — \textquotedblleft{}to get down, to get off\textquotedblright{}.

\begin{grammarlens}[title={What you've built}]
One more word for this chapter.
\end{grammarlens}

\subsection*{Your turn}
\begin{itemize}
\raggedright
  \item \emph{Say it:} \emph{digu}
  \item \emph{Say it:} \emph{digu}, once more
  \item \emph{Say it:} say \emph{digu}
  \item \emph{Recall:} say the Telugu for to forget, then the Telugu for to jump, then say \emph{digu} again
\end{itemize}

\begin{tcolorbox}[breakable,colback=teal!4,colframe=teal!35!black,title={Before you move on}]
What does \te{దిగు} mean? (\textquotedblleft{}To get down, to get off\textquotedblright{}.) Say it once more.
\end{tcolorbox}
